%----------------------------------------------------------------------------------------
% NOTATIE EN FIGUREN VAN DE CURSUS COMPLEXE GETALLEN
% Staat in de preambule en niet in een hoofdstuk: ook de lijsten met oplossingen gebruiken dit.
%----------------------------------------------------------------------------------------
% Bewaarde grafiek, geplaatst door maak-pdf.py: bestand en pagina in figuren-cache/
\newcommand{\grafiekbestand}[2]{\includegraphics[page=#2,trim=3bp 3bp 3bp 3bp]{figuren-cache/#1}}

% Getallenverzamelingen, reëel en imaginair deel, toegevoegde
\newcommand{\C}{\mathbb{C}}
\newcommand{\R}{\mathbb{R}}
\renewcommand{\Re}{\operatorname{Re}}
\renewcommand{\Im}{\operatorname{Im}}
\newcommand{\conj}[1]{\overline{#1}}
% Vlaamse notatie voor de inverse tangens, zoals in Integralen
\newcommand{\bgtan}{\operatorname{bgtan}}

% Zoals stappen, maar met een bredere kolom voor de berekening: voor lange formules
\newenvironment{stappenbreed}{%
  \par\renewcommand{\arraystretch}{1.25}%
  \noindent\begin{tabular}{@{}>{$\displaystyle}p{9.8cm}<{$}@{\hspace{4mm}}>{\RaggedRight\small\itshape\arraybackslash}p{4.9cm}@{}}%
}{\end{tabular}\par}

% Zoals stappen, met eigen kolombreedtes: \begin{stappenvrij}{6.2cm}{3.8cm}. Voor naast een figuur.
\newenvironment{stappenvrij}[2]{%
  \par\renewcommand{\arraystretch}{1.25}%
  \noindent\begin{tabular}{@{}>{$\displaystyle}p{#1}<{$}@{\hspace{3mm}}>{\RaggedRight\small\itshape\arraybackslash}p{#2}@{}}%
}{\end{tabular}\par}

% Tekst of berekening links, figuur rechts: \naast[breedte rechts]{links}{rechts}
% De breedte rechts komt uit de figuur zelf; het optionele argument wordt niet gebruikt.
\newsavebox{\naastbox}
\newcommand{\naast}[3][]{%
  \par\noindent\sbox{\naastbox}{#3}%
  \begin{minipage}[t]{\dimexpr\linewidth-\wd\naastbox-3mm\relax}\vspace{0pt}#2\end{minipage}\hfill
  \begin{minipage}[t]{\wd\naastbox}\vspace{0pt}\usebox{\naastbox}\end{minipage}\par\addvspace{1mm}}
% Een rij figuren, gecentreerd
\newcommand{\figuren}[1]{\par\addvspace{1mm}{\centering #1\par}\addvspace{1.5mm}}
% Opgave met een grafiek rechts: \begin{metgrafiek}[breedte]{tikz} opgave en deelvragen \end{metgrafiek}
% De solution komt na \end{metgrafiek}, zoals bij metprent.
\newsavebox{\grafiekbox}
\newenvironment{metgrafiek}[2][]{%
  \sbox{\grafiekbox}{#2}%
  \begin{minipage}[t]{\dimexpr\linewidth-\wd\grafiekbox-3mm\relax}%
}{%
  \end{minipage}\hfill
  \begin{minipage}[t]{\wd\grafiekbox}\vspace{-2mm}\usebox{\grafiekbox}\end{minipage}\par}

% HET COMPLEXE VLAK (vlak van Gauss)
% \gaussleeg{xmin}{xmax}{ymin}{ymax}: rooster en assen Re en Im, zonder getallen.
% \gauss{xmin}{xmax}{ymin}{ymax}: hetzelfde, met de gehele getallen op de assen (1, 2, ... en i, 2i, ...).
% \gaussstap{stap}{xmin}{xmax}{ymin}{ymax}: zoals \gauss, met een getal bij elke stap-de lijn.
\newcommand{\gaussleeg}[4]{%
  \draw[black!14,very thin] (#1,#3) grid (#2,#4);
  \draw[->,grijs] (#1,0) -- ({#2+.45},0) node[right,text=black,font=\scriptsize] {Re};
  \draw[->,grijs] (0,#3) -- (0,{#4+.45}) node[above,text=black,font=\scriptsize] {Im};}
\newcommand{\gaussstap}[5]{%
  \gaussleeg{#2}{#3}{#4}{#5}
  \foreach \t in {#2,...,#3}{\pgfmathtruncatemacro{\gr}{mod(abs(\t),#1)}%
    \ifnum\t=0\else\ifnum\gr=0\node[below,inner sep=1.5pt,text=black,font=\scriptsize] at (\t,0) {$\t$};\fi\fi}
  \foreach \t in {#4,...,#5}{\pgfmathtruncatemacro{\gr}{mod(abs(\t),#1)}%
    \ifnum\t=0\else\ifnum\gr=0\node[left,inner sep=1.5pt,text=black,font=\scriptsize] at (0,\t)
      {$\ifnum\t=1 i\else\ifnum\t=-1 -i\else\t i\fi\fi$};\fi\fi}}
\newcommand{\gauss}[4]{\gaussstap{1}{#1}{#2}{#3}{#4}}
% Een complex getal als punt met een label: \cpunt[positie]{x}{y}{label}
\newcommand{\cpunt}[4][above right]{\fill[accent] (#2,#3) circle (2.2pt) node[#1,inner sep=1.5pt,text=black,font=\small] {#4};}
\tikzset{
  % Een complex getal als vector vanuit de oorsprong: \draw[vector] (0,0) -- (3,2);
  vector/.style={-{Stealth[length=2.2mm]},accent,thick},
  hulplijn/.style={grijs,dashed,thin},
}
